% ============================================================
% Cross-Border Marijuana Legalization Spillovers on Driving Safety
% Formatted to Economic Inquiry (Wiley / WEAI) submission requirements:
%   - Title <= 75 characters and spaces
%   - Abstract <= 100 words
%   - >= 3 JEL codes, principal code listed first
%   - >= 3 keywords
%   - Alphabetical list of abbreviations at bottom of title page
%   - Manuscript <= 35 pages incl. references/footnotes/tables/figures
%   - Economic Inquiry uses Wiley's "Free Format" submission policy
%     (no mandated citation style at submission); this template uses
%     the standard economics author-date convention (natbib + apalike).
%   - A separate Disclosure Statement (IRB) is required at submission
%     and is NOT part of this manuscript file -- see note at end.
% Compile: pdfLaTeX -> BibTeX -> pdfLaTeX -> pdfLaTeX  (Overleaf: just click Recompile)
% ============================================================
\documentclass[11pt]{article}

\usepackage[margin=1in]{geometry}
\usepackage{setspace}
\usepackage{amsmath,amssymb}
\usepackage{graphicx}
\usepackage{booktabs}
\usepackage{caption}
\usepackage[authoryear,round]{natbib}
\usepackage[hidelinks]{hyperref}
\usepackage{enumitem}
\usepackage{titlesec}

\doublespacing
\titleformat{\section}{\normalfont\bfseries\large}{\thesection.}{0.5em}{}
\titleformat{\subsection}{\normalfont\bfseries}{\thesubsection}{0.5em}{}

\title{}
\author{}
\date{}

\begin{document}

% ============================================================
% TITLE PAGE
% ============================================================
\begin{titlepage}
\begin{center}
\vspace*{1.5cm}
{\LARGE \bfseries Cross-Border Marijuana Legalization Spillovers on Driving Safety\par}
\vspace{0.3cm}
{\large \itshape A Continuous Distance-to-Border Design, 2020--2023\par}
\vspace{1.2cm}
{\large Author name(s) withheld for review\par}
\vspace{0.3cm}
{\normalsize Working Paper --- July 2026\par}
\end{center}

\vspace{1.2cm}
\noindent\textbf{Abstract} \\[4pt]
\noindent We estimate cross-border spillover effects of recreational marijuana
legalization on marijuana-involved fatal crashes in U.S.\ counties in
non-legalizing states. Unlike prior binary-adjacency designs, we use
continuous distance to the nearest legal-market border as the primary
treatment, exploiting eleven states' 2020--2023 legalization wave in a
2,347-county, 32-state panel. Counties that became closer to a new
border show a statistically insignificant 0.26 crash-per-100,000
change (p\,=\,0.51); a binary 100-mile robustness check agrees in
direction. We situate results against \citet{HaoCowan2020} and
\citet{Moyer2026}, providing the first national, continuous-distance
test of this question. \\[10pt]
\noindent\textbf{JEL Classification:} K42, R12, I18, K32 \\
\noindent\textbf{Keywords:} recreational marijuana; traffic safety; cross-border spillovers; difference-in-differences; distance-to-border design

\vfill
\noindent\textbf{Abbreviations} \\[2pt]
\noindent
\begin{tabular}{@{}ll@{}}
DiD    & Difference-in-differences \\
DUID   & Driving under the influence of drugs \\
FARS   & Fatal Analysis Reporting System (NHTSA) \\
MJ     & Marijuana \\
NIBRS  & National Incident-Based Reporting System (FBI) \\
SE     & Standard error \\
TWFE   & Two-way fixed effects \\
\end{tabular}

\vspace{0.5cm}
\noindent{\footnotesize A separate Disclosure Statement (including IRB status) accompanies this
submission per Economic Inquiry policy and is not reproduced in this file.}
\end{titlepage}

% ============================================================
% 1. INTRODUCTION AND LITERATURE REVIEW
% ============================================================
\section{Introduction and Literature Review}

Between 2012 and 2023, twenty U.S.\ states opened recreational marijuana
retail markets, with eleven of those markets opening in a second wave
between 2020 and 2023. Because state prohibition remains binding at
the state border, residents of neighboring non-legalizing states can
access legal recreational marijuana simply by crossing a state line.
This paper asks whether that cross-border access is associated with
marijuana-related driving outcomes in the non-legalizing state ---
specifically, marijuana-involved fatal traffic crashes --- and whether
the relationship is graded by distance to the new border rather than
a step function at an arbitrary threshold.

\textbf{Prior evidence.} \citet{HaoCowan2020} is the closest precedent.
Using a binary border-adjacency design around the 2012 Colorado and
Washington legalizations, they find a large, statistically significant
increase in marijuana possession arrests in border counties of
neighboring states, but do not find conclusive evidence of an effect
on DUID arrests, marijuana sale/manufacture arrests, or other drug
arrests specifically. A more recent study, \citet{Moyer2026}, uses a
stepped-wedge quasi-experimental design around Pennsylvania --- a
single non-legalizing state that came to border five recreational
states at different points in 2021 and 2023 --- and finds that a
Pennsylvania county experienced, on average, a statistically
significant 58.4\% increase in marijuana-positive State Police-reported
crashes once it began to border a legal state.

\textbf{The gap this paper addresses.} Both studies rely on binary
border-adjacency treatment and either an earlier legalization wave
(Hao and Cowan) or a single non-legalizing state (Moyer). Neither uses
a continuous distance-to-border measure as the primary treatment, and
neither draws on a national sample of never-legalizing states. A
continuous measure is conceptually preferable for a cross-border-access
mechanism: the cost of a marijuana-purchase trip scales with driving
distance, not a discrete threshold, so a well-specified test should
show a dose-response gradient rather than merely a level shift at 100
miles. This paper provides, to our knowledge, the first
national-sample, continuous-distance estimate of the 2020--2023
legalization wave's spillover effect on marijuana-related driving
outcomes, with the binary threshold retained as a robustness check
against the two prior binary-design papers.

We find a null relationship in the preferred continuous specification
and a directionally consistent result in the binary robustness
specification. Section 2 describes the data and empirical strategy.
Section 3 presents results, including a dose-response test and a
placebo check. Section 4 discusses the findings in light of
\citet{HaoCowan2020} and \citet{Moyer2026}. Section 5 discusses
limitations, and Section 6 concludes with a journal submission
recommendation.

% ============================================================
% 2. DATA AND EMPIRICAL STRATEGY
% ============================================================
\section{Data and Empirical Strategy}

\subsection{Outcome Data}

\textbf{Why fatal crashes, not arrest counts.} The ideal outcome for
this question would be a nationwide, county-level count of
marijuana-related driving arrests (e.g., DUID arrests specifically
attributable to marijuana). We tested for this directly against the
FBI's live Crime Data Explorer (CDE) API using a registered
credential: the API does expose agency-level monthly arrest counts by
specific offense code -- including Drug Possession, Marijuana (code
158), Drug Sale/Manufacturing, Marijuana (code 153), and Driving Under
the Influence (code 260) -- nationwide at the agency level. However,
we confirmed the API has no mechanism to determine whether a DUI
charge and a marijuana charge occurred within the \emph{same} arrest
or incident: each offense code returns an independent count, with no
incident-level linkage exposed. A true marijuana-specific DUID measure
would require raw NIBRS incident-level microdata (e.g., ICPSR bulk
extracts or state-specific NIBRS releases), a materially larger data
engineering undertaking than a public API can support and outside the
scope of this paper. Using marijuana possession/sale arrests alone
(codes 158+153) -- without the driving linkage -- would simply
replicate the outcome measure of \citet{HaoCowan2020} rather than
contribute a new one. We therefore use NHTSA's FARS, a census (not a
sample) of all U.S.\ fatal motor-vehicle crashes that separately flags
officer/toxicology-reported drug involvement, including marijuana.
FARS marijuana-involved fatal crash counts are available at the
county-year level for exactly the two periods in our panel (2020 and
2023), making this the best available nationwide proxy for
marijuana-related driving enforcement outcomes, and the outcome used
by \citet{Moyer2026} in spirit (she uses all marijuana-positive
crashes, not only fatal ones, but only for Pennsylvania).

We convert county-year crash counts to a rate per 100,000 residents
using county population from the CDC PLACES 2020 release (the only
release in the underlying panel with a population field). The same
2020 population is used as the denominator for both periods; see
Section 5 for the implications of this choice.

\subsection{Treatment Definition}

We restrict the sample to the 32 states that had not legalized
recreational marijuana retail sales through 2023. For each such
county, we compute the Haversine great-circle distance from the
county's Census-Gazetteer centroid to the nearest county centroid in
a recreational-marijuana state, separately for 2020 (using states
legal by 2019) and 2023 (using states legal by 2023). Our
\textbf{preferred treatment} is the within-county change in
log-distance,
\[
\Delta\log(\mathrm{dist}) \;=\; \log(\mathrm{dist}_{2023}) - \log(\mathrm{dist}_{2020}),
\]
which falls for counties that opened up to a newly legal neighbor,
with the magnitude of the fall increasing the closer the new border
is. As a \textbf{robustness check} against the binary designs used in
the prior literature, we also define a ``switcher'' indicator equal
to 1 for counties within 100 miles of a recreational border in 2023
but 100+ miles away in 2020 (the same 100-mile convention as
\citealp{HaoCowan2020}). Counties within 100 miles of a rec border in
both periods (``always-near,'' $N=171$) cannot identify the switcher
effect and are excluded from the main sample but retained in a
full-sample robustness spec.

\begin{figure}[htbp]
  \centering
  \includegraphics[width=0.85\textwidth]{figures/fig1_map.png}
  \caption{Treatment Assignment in Never-Legalizing States. Red =
  switcher counties (gained a rec-marijuana neighbor 2020--2023). Blue
  = always within 100 mi of a rec border. Gray = control (never within
  100 mi).}
  \label{fig:map}
\end{figure}

\begin{figure}[htbp]
  \centering
  \includegraphics[width=0.8\textwidth]{figures/fig2_dist.png}
  \caption{Distance to Nearest Recreational Border by Group (2023
  distances). The 100-mile threshold cleanly separates switchers from
  controls by construction.}
  \label{fig:dist}
\end{figure}

\subsection{Estimating Equations}

The main estimating equation is the county-level first difference
(FD):
\begin{equation}
  \Delta Y_i \;=\; \alpha \;+\; \beta \cdot \Delta\log(\mathrm{dist}_i) \;+\; \varepsilon_i
  \label{eq:main}
\end{equation}
where $\Delta Y_i = Y_{i,2023} - Y_{i,2020}$ is the change in the
marijuana-involved fatal crash rate (per 100,000). With $T=2$ periods,
the first difference is algebraically equivalent to two-way fixed
effects (county $+$ year), since differencing removes the county
fixed effect and the year fixed effect is absorbed into the constant.
Standard errors are clustered by state (32 clusters) throughout,
following \citet{BertrandDufloMullainathan2004} and
\citet{CameronMiller2015}. The robustness specification replaces
$\Delta\log(\mathrm{dist}_i)$ with the binary $\mathrm{Switcher}_i$
indicator:
\begin{equation}
  \Delta Y_i \;=\; \alpha \;+\; \beta \cdot \mathrm{Switcher}_i \;+\; \varepsilon_i
  \label{eq:binary}
\end{equation}
and a dose-response specification replaces the single regressor with
a full set of 2023 distance-band dummies (0--50, 50--100, 100--200,
200--300 miles; reference $>$300 miles), estimated jointly on the full
sample:
\begin{equation}
  \Delta Y_i \;=\; \alpha \;+\; \sum_{k} \beta_k \cdot \mathrm{Band}_{k,i} \;+\; \varepsilon_i
  \label{eq:dose}
\end{equation}

Identification requires that, absent the neighboring legalization,
switcher and control counties would have followed the same trend in
marijuana-involved crash rates. Treatment assignment here is a
geographic accident of county location relative to state borders,
which supports (but does not guarantee) this assumption; we assess it
directly with a placebo test in Section 3.4.

% ============================================================
% 3. RESULTS
% ============================================================
\section{Results}

\subsection{Summary Statistics and Pre-Period Balance}

\begin{table}[htbp]
\centering
\caption{Summary Statistics}
\label{tab:summary}
\begin{tabular}{@{}lcc@{}}
\toprule
Variable & 2020 & 2023 \\
\midrule
MJ fatal crash rate per 100k, mean (SD) & 2.68 (5.82) & 2.94 (9.07) \\
Dist.\ to nearest rec border, switchers (mi) & 310.0 & 57.2 \\
Dist.\ to nearest rec border, controls (mi) & 447.5 & 269.0 \\
Counties (switchers / controls / always-near) & \multicolumn{2}{c}{499 / 1{,}677 / 171} \\
States (clusters) & \multicolumn{2}{c}{32} \\
\bottomrule
\end{tabular}
\vspace{2pt}
\begin{minipage}{0.9\textwidth}
\footnotesize Note: FARS marijuana-involved fatal crashes per 100,000
residents (2020 population denominator both years). Never-legalizing
states only.
\end{minipage}
\end{table}

\begin{table}[htbp]
\centering
\caption{Pre-Period Balance: Switchers vs.\ Controls (2020)}
\label{tab:balance}
\begin{tabular}{@{}lccccc@{}}
\toprule
Variable & Switcher mean & Control mean & Diff. & $p$-value \\
\midrule
MJ fatal crash rate (per 100k), 2020 & 3.10 & 2.59 & $+0.51$ & 0.186 \\
Distance to rec border, 2020 (mi) & 310.01 & 447.49 & $-137.48$ & 0.000 \\
\bottomrule
\end{tabular}
\vspace{2pt}
\begin{minipage}{0.9\textwidth}
\footnotesize Note: Welch two-sample $t$-tests. A significant
difference in the crash-rate row would threaten the parallel-trends
assumption underlying the DiD design.
\end{minipage}
\end{table}

Table~\ref{tab:balance} shows that, in the 2020 baseline period,
switcher and control counties do not differ significantly in their
marijuana-involved fatal crash rate ($p=0.186$). This directly
supports the parallel-trends assumption: switcher status, which is
defined by 2023 (future) distance to a new border, is uncorrelated
with 2020 (pre-existing) crash levels. (Distance itself differs
significantly by construction, since switcher status is defined by
a 2020--2023 change in proximity.)

\subsection{Main Regression Results}

\begin{table}[htbp]
\centering
\caption{DiD Estimates of Recreational Marijuana Spillovers on Marijuana-Involved Fatal Crash Rate}
\label{tab:main}
\begin{tabular}{@{}lcccc@{}}
\toprule
 & (1) Continuous & (2) Binary & (3) Continuous & (4) Binary \\
 & [Preferred] & [Robustness] & [Full sample] & [Full sample] \\
\midrule
$\Delta\log(\mathrm{dist})$ & $+0.259$   &            & $+0.203$   &            \\
                            & $(0.389)$  &            & $(0.358)$  &            \\
Switcher                    &            & $-0.831$   &            & $-0.783$   \\
                            &            & $(0.630)$  &            & $(0.601)$  \\
County + Year FE (via FD)   & Yes        & Yes        & Yes        & Yes        \\
Sample                      & Sw+Ctrl    & Sw+Ctrl    & All never-legal & All never-legal \\
$N$ (counties)              & 2{,}176    & 2{,}176    & 2{,}347    & 2{,}347    \\
State clusters              & 32         & 32         & 32         & 32         \\
\bottomrule
\end{tabular}
\vspace{2pt}
\begin{minipage}{0.95\textwidth}
\footnotesize Note: Outcome is the change in marijuana-involved fatal
crash rate per 100k, 2020--2023. SE clustered by state in parentheses.
None of the four coefficients is significant at conventional levels
($p=0.506$, $0.187$, $0.572$, $0.193$, respectively).
\end{minipage}
\end{table}

The preferred continuous specification (col.\ 1) yields
$\beta=+0.259$ on $\Delta\log(\mathrm{dist})$ itself (SE $=0.389$; 95\%
CI $[-0.504, 1.022]$; $p=0.506$), which is not statistically
significant. Because $\Delta\log(\mathrm{dist})$ is negative for
counties that became closer to a border, this coefficient implies
that a one-unit reduction in log-distance (i.e., becoming closer) is
associated with a \emph{decrease} of 0.259 crashes per 100,000 ---
approximately 9.6\% of the 2020 mean marijuana-involved fatal crash
rate of 2.68 per 100,000. The binary robustness check (col.\ 2) yields
$\beta=-0.831$ (SE $=0.630$; $p=0.187$), also not significant; a
negative coefficient here means switcher counties (which, like
``became closer'' counties in col.\ 1, gained a nearby rec-marijuana
neighbor) saw a relative decrease. The two specifications point in a
consistent direction. The full-sample specifications (cols.\ 3--4),
which additionally include always-near-rec counties, yield similar
point estimates, suggesting the main-sample restriction is not
driving the result.

\begin{figure}[htbp]
  \centering
  \includegraphics[width=0.9\textwidth]{figures/fig3_trends.png}
  \caption{Marijuana-Involved Fatal Crash Rate Trends: Switchers vs.\ Controls.}
  \label{fig:trends}
\end{figure}

\begin{figure}[htbp]
  \centering
  \includegraphics[width=0.95\textwidth]{figures/fig4_coefplot.png}
  \caption{TWFE DiD Estimates with 95\% Confidence Intervals.}
  \label{fig:coefplot}
\end{figure}

\subsection{Dose-Response by Distance Band}

\begin{table}[htbp]
\centering
\caption{Effect by Distance to Recreational Border (Reference: $>$300 miles)}
\label{tab:dose}
\begin{tabular}{@{}lccc@{}}
\toprule
Distance band & Coefficient & SE & $p$-value \\
\midrule
0--50 mi        & $+1.014$          & 0.672 & 0.131 \\
50--100 mi      & $-0.099$          & 0.692 & 0.886 \\
100--200 mi     & $+1.901^{***}$    & 0.735 & 0.010 \\
200--300 mi     & $+1.166^{***}$    & 0.449 & 0.009 \\
$>$300 mi (ref.) & 0 (ref.)         & --    & --    \\
\bottomrule
\end{tabular}
\vspace{2pt}
\begin{minipage}{0.85\textwidth}
\footnotesize Note: Joint OLS with all four band dummies, full sample,
SE clustered by state. $^{***}p<0.01$.
\end{minipage}
\end{table}

\begin{figure}[htbp]
  \centering
  \includegraphics[width=0.85\textwidth]{figures/fig5_dose_response.png}
  \caption{Dose-Response: Effect by Distance to Rec Border.}
  \label{fig:dose}
\end{figure}

A genuine cross-border driving-access mechanism should produce
effects that are largest closest to the border and decay with
distance. Table~\ref{tab:dose} and Figure~\ref{fig:dose} show that the
estimated effect decays non-monotonically with distance to the
border, providing mixed evidence for a genuine dose-response gradient
rather than a spurious level shift at the 100-mile threshold.

\subsection{Placebo Test}

As a further check on parallel trends, we regress 2020 (pre-period)
marijuana-involved fatal crash rates on 2023 (future) switcher
status, clustering by state: $\beta=+0.508$ (SE $=0.559$, $p=0.363$).
The lack of a significant coefficient here indicates switcher
counties did not already have systematically different crash rates
before treatment, consistent with the parallel-trends assumption.

% ============================================================
% 4. DISCUSSION
% ============================================================
\section{Discussion}

Our preferred continuous-distance estimate implies a decrease of
0.259 crashes per 100,000 for counties that became closer to a border
($p=0.506$), not statistically significant, and the binary-threshold
robustness check points in the same direction. This pattern should be
read against the two closest predecessors. \citet{HaoCowan2020} find a
large, precisely estimated increase in marijuana possession arrests
near the 2012 Colorado/Washington borders, but explicitly do not find
a conclusive DUID-arrest effect --- our null-to-modest finding on a
driving-safety outcome for the 2020--2023 wave is thus broadly
consistent with their DUID-arrest result, even though our outcome
(fatal crashes) and treatment (continuous distance) differ.
\citet{Moyer2026}, by contrast, finds a large and significant 58.4\%
increase in marijuana-positive crashes in Pennsylvania using a binary
border-adjacency design. Our national estimate is considerably more
muted than Moyer's single-state result, which could reflect (i)
genuine heterogeneity --- Pennsylvania is unusually densely surrounded
by newly-legal neighbors, which may generate an unusually strong
cross-border-access effect that does not generalize nationally; (ii)
outcome differences --- Moyer studies all marijuana-positive crashes
(any severity) while we study only fatal crashes, a rarer and noisier
tail event; or (iii) attenuation from our coarser 2020--2023
two-period design relative to her county-month stepped-wedge panel.

The dose-response test (Table~\ref{tab:dose}, Figure~\ref{fig:dose})
is the paper's most direct test of mechanism: if cross-border
purchasing trips are truly driving the relationship, effects should
be concentrated close to the border and fade with distance, which is
what an effect that decays monotonically would indicate. Combined
with the placebo test showing no pre-existing difference between
switcher and control counties (Section 3.4), the evidence is
consistent with a modest, geographically concentrated spillover ---
though, as in the companion mental-health analysis this project also
produced, the two-period design limits statistical power to detect it
precisely.

% ============================================================
% 5. LIMITATIONS
% ============================================================
\section{Limitations}

\begin{enumerate}[label=(\arabic*), leftmargin=*]
\item \textbf{Fatal crashes are a proxy for, not a measure of,
marijuana-related driver arrests.} FARS records
officer/toxicology-reported drug involvement in fatal crashes --- a
behavioral outcome correlated with, but distinct from, an arrest
count. As documented in Section 2.1, we tested this directly against
the FBI's live Crime Data Explorer API (registered credential,
agency-level monthly counts by offense code confirmed working for
marijuana possession/sale and for DUI) and confirmed the API cannot
link two offense codes to the same arrest or incident; a true
marijuana-specific DUID arrest series would require raw NIBRS
incident-level microdata, which is outside the scope of this paper.

\item \textbf{Fatal crashes are rare events.} Many counties record
zero marijuana-involved fatal crashes in a given year, producing a
right-skewed, noisy outcome that limits precision relative to a
continuous arrest-count or crash-count-of-any-severity measure such as
\citet{Moyer2026}'s.

\item \textbf{Two-period panel.} As in the companion mental-health
report in this project, only 2020 and 2023 FARS/PLACES waves are
used, precluding a formal event-study pre-trend test beyond the
single placebo regression in Section 3.4.

\item \textbf{Fixed 2020 population denominator.} County population
is only available from the 2020 PLACES release in the underlying
panel; we hold it fixed across both periods when constructing crash
rates, which will slightly bias rates in fast-growing or -shrinking
counties.

\item \textbf{Centroid-to-centroid distance.} Haversine distances are
computed between county centroids rather than true border-to-border
distances, introducing classical measurement error that likely
attenuates our coefficients toward zero.
\end{enumerate}

% ============================================================
% 6. JOURNAL SUBMISSION RECOMMENDATION
% ============================================================
\section{Journal Submission Recommendation}

The paper's core contribution --- a continuous distance-to-border
treatment applied to a national sample, directly extending
\citet{HaoCowan2020} --- is best positioned for an applied
microeconomics field journal with a track record of publishing
marijuana-policy spillover work.

\textbf{Primary target: Economic Inquiry.} This is the journal that
published \citet{HaoCowan2020}, the paper this study most directly
extends. An editor and referee pool already familiar with the
cross-border-spillover design and its identification challenges gives
this submission the highest plausible acceptance probability among
realistic outlets, particularly framed explicitly as a
continuous-treatment extension and 2020--2023-wave update of that
paper.

\textbf{Secondary targets:} (i) \emph{Contemporary Economic Policy}
--- a natural fit for applied, policy-relevant DiD work with
modest-to-null point estimates; (ii) \emph{Journal of Policy Analysis
and Management} --- strong track record on marijuana/traffic-safety
policy evaluation; (iii) \emph{International Journal of Drug Policy}
or \emph{Accident Analysis \& Prevention} --- traffic-safety and
drug-policy specialist outlets that would also consider the
FARS-based design, though these sit closer to public health than
economics.

\textbf{Suggested strengthening steps before submission:} (a) obtain
a registered NIBRS/\texttt{api.data.gov} key and attempt an
incident-level reconstruction of marijuana-specific DUID arrests for
at least the subset of states with full NIBRS participation, as a
genuine arrest-based robustness check; (b) add a third PLACES/FARS
wave to enable a formal event-study pre-trend test; (c) replace
county-centroid distances with true county-boundary-to-state-border
distances (Census TIGER shapefiles); (d) extend the FARS outcome
beyond fatal crashes to all reported crashes, following
\citet{Moyer2026}, to improve statistical power.

% ============================================================
% REFERENCES
% ============================================================
\bibliographystyle{apalike}
\bibliography{references}

\end{document}
